% 2011 年山东卷文科第 14 题：程序框图。
% 来源 PDF：普通高考/2011/2011山东文.pdf，SHA-256 bf61cf7bb39a3ccfc57735cd2a7ab34387cbfa6bc86a5e7f8246e13292b0e023。
% 原页：第 1 页（A3 三栏）第 3 栏上部，书页码 691。
% 重绘依据：按原页矢量层逐框实测复刻。坐标取原页 pt：X = x_页 − 957.56，Y = 302.20 − y_页。
%   8 个框：开始/结束为圆角矩形 27.98×16.83、27.98×17.00（圆角半径 3.98）；
%   输入、输出为平行四边形（上下边水平、左右边同向斜，横向偏移 6.23，高 17.13）；
%   两个判断为菱形（半对角线 64.32×16.02 与 42.63×10.59）；
%   三个语句框为矩形：y=105 为 42.25×15.80，y=70l+21m+15n 为 102.52×16.33，y=y−105 为 60.61×15.80。
%   连线：主干 6 段竖直，另有 4 段折线——「是」支出菱形左顶点左行 22.63 后下行进 y=105 框；
%   y=105 框底下行 63.30 再右行汇入主干（在菱形 2 与输出之间）；「是」支出菱形右顶点右行 53.50
%   后上行进 y=y−105 框；该框顶上行 5.87 再左行汇入主干（在菱形 2 之上）。全图无虚线、无圆点。
\begin{tikzpicture}[x=0.035cm, y=0.035cm, >=stealth, semithick, font=\small,
  every node/.style={inner sep=1.5pt}]

  % ===== 8 个框 =====
  \draw[rounded corners=3.98] (-13.99,228.40) rectangle (13.99,245.23);                    % 开始
  \draw (-58.82,196.69) -- (-52.59,213.82) -- (58.80,213.82) -- (52.56,196.69) -- cycle;     % 输入
  \draw (0,182.04) -- (64.32,166.02) -- (0,150.00) -- (-64.32,166.02) -- cycle;              % 菱形 1
  \draw (-108.08,119.29) rectangle (-65.83,135.09);                                          % y=105
  \draw (-51.26,119.03) rectangle (51.26,135.36);                                            % y=70l+21m+15n
  \draw (0,84.54) -- (42.63,73.95) -- (0,63.36) -- (-42.63,73.95) -- cycle;                  % 菱形 2
  \draw (65.82,88.66) rectangle (126.43,104.46);                                             % y=y-105
  \draw (-24.74,31.57) -- (-18.51,48.70) -- (24.73,48.70) -- (18.50,31.57) -- cycle;         % 输出
  \draw[rounded corners=3.98] (-13.99,0) rectangle (13.99,17.00);                            % 结束

  % ===== 主干（竖直） =====
  \draw[->] (0,228.40) -- (0,213.82);            % 开始 → 输入
  \draw[->] (0,196.69) -- (0,182.04);            % 输入 → 菱形 1
  \draw[->] (0,150.00) -- (0,135.36);            % 菱形 1（否）→ y=70l+21m+15n
  \draw[->] (0,119.03) -- (0,84.54);             % y=70l+21m+15n → 菱形 2
  \draw[->] (0,63.36) -- (0,48.70);              % 菱形 2（否）→ 输出
  \draw[->] (0,31.57) -- (0,17.00);              % 输出 → 结束

  % ===== 「是」支与回流折线 =====
  \draw[->] (-64.32,166.02) -- (-86.95,166.02) -- (-86.95,135.09);   % 是 → y=105
  \draw[->] (-86.95,119.29) -- (-86.95,55.99) -- (0,55.99);          % y=105 → 主干（输出前）
  \draw[->] (42.63,73.95) -- (96.13,73.95) -- (96.13,88.66);         % 是 → y=y-105
  \draw[->] (96.13,104.46) -- (96.13,110.33) -- (0,110.33);          % y=y-105 → 主干（菱形 2 前）

  % ===== 框内文字 =====
  \node at (0.34,236.85) {开始};
  \node at (0.3,205.75) {输入非负整数 $l$，$m$，$n$};
  \node at (0,165.60) {$l^2+m^2+n^2=0$};
  \node at (-86.61,127.50) {$y=105$};
  \node at (0.34,127.40) {$y=70l+21m+15n$};
  \node at (0.39,74.25) {$y>105$};
  \node at (96.54,96.85) {$y=y-105$};
  \node at (0.19,40.60) {输出 $y$};
  \node at (0.34,9.15) {结束};

  % ===== 分支标签 =====
  \node at (-70.5,175.2) {是};
  \node at (9.3,143.4) {否};
  \node at (43.3,83.2) {是};
  \node at (9.3,56.7) {否};

\end{tikzpicture}
